\section{程序组织与执行顺序}
\label{sec:program}
\code{config.py} 定义物性、几何和单位换算；\code{lattice.py} 定义 D2Q9 原语；
\code{simulator.py} 包含流体、刚体、热更新；\code{reporting.py} 和示例处理输出。
所有入口和测试均导入 \code{iceflow2d_v3}。

\subsection{完整时间步}
\begin{lstlisting}[language=Python]
for _ in range(num_steps):
    self._advance_fluid()
    self._advance_body()
    self._advance_thermal()
    self.steps += 1
\end{lstlisting}
流体阶段包括两组碰撞、流送、宏观量与压力恢复，以及局部相场限幅。
刚体阶段包括积分、墙面约束、光栅化、掩码更新和运动释放格点重建。
热阶段每步重映射水温并对流水显热；累计满热更新间隔后执行导热、界面换热和材料删除。
随后更新剩余冰体质量性质和锐界，把融化释放格点置为空气，并同步热水开口。

\subsection{初始化与采样}
先初始化材料网格、刚体位姿、锐界和气液分布，再进行相场预热及局部限幅。
最后初始化水显热和冻结静水参考场。采样时刻对齐慢热更新边界。
\begin{lstlisting}[language=bash]
python -m iceflow2d_v3 --end-time-s 0.00005 \
  --output-interval-s 0.00005 --no-plot --save-npz
python -m pytest iceflow2d_v3 -q
\end{lstlisting}

\subsection{结果读取}
CSV、NPZ 和元数据保存位姿、速度、剩余质量、惯量、融化比例、当前水量和平均水温。
不保存全局守恒残差、液水目标或由融化产生的流体源项统计。
体积焓是当前冰显热与当前液水参考焓的显示量，不累计已经消失材料的潜热。
PNG/GIF 按帧写出；仅在请求 NPZ 时保留所有采样数组。
